\scanpage{10}
\handwritten{24.4\unsure{,|.}67}\\

Dear Murre,\\

I am sending you \unsure{enchosed|enclosed} the sketch which I promised on ramification data, as well as your manuscript together with my annotations. As the paragraphs 1 to 5 will have to be changed entirely, there is maybe not much point for you to look up to carefully my annotations on these. \deleted{As} I would only suggest that you put more \unsure{formmulas|formulas}, arrows etc on \deleted{\unsure{separate}} separate lines, for easier reading; that you normalize the \unsure{stadfard|stanfard|standard} references to SGA, EGA (for instance, no need to indicate SGA when reference is being made to the seminar 1960-61 in which your talk is going to be included itself, reference as (XI 6.1) is then sufficient); that you be cautious when using a locution like ``if primitive $n$.th roots exist over $S$'', which does not mean much by itself, unless you make clear what you really mean. I would propose the following tentative plan for your talk:

\begin{enumerate}[label={}]
\item[1.] Kummer coverings with respect to a given set of sections.
\item[2.] Ramification data.
\item[3.] Properties of the catégory of coverings of type $\underline{R}$. The group ${}_{1}^{\underline{R}}(S,\ )$
% No base glyph is visible before the stacked R and 1; the second argument is blank.
\item[4.] Functorial behaviours.
\item[5.] Galois coverings of type $\underline{R}$.
\item[6.] Abhyankars theorem (for divisors with normal crossings).
\end{enumerate}
etc

The sections 6 to 10 essentially the same as in your notes, just taking into account my observations, and changing inner references according to the changes made in 1. to 5. A paragraph 11 could be an appendix on the tame variant of Lefschetz's theorem. - Par.1 is essentially your previous par.1 ; you should check exactly what is needed in it for \unsure{fhat|what|that}
% archive sheet number: 9
